\documentclass[10pt,a4paper]{amsart}
\usepackage{amsmath,amssymb,amsthm,graphicx}
% Retained preamble for the Stacks Project chapter snapshot in this directory.
% The chapter's own preamble.tex is a copy of the shared Stacks Project preamble; the block
% below is sliced MECHANICALLY, line for line, from it. Nothing was added, removed or
% redefined; the wrapper class, the xr-hyper/hyperref wiring and the externaldocument list
% are omitted because the assembler supplies the document class.
%
% lines taken from preamble.tex: 3, 5, 9, 11, 12, 156, 157, 158, ... (51 lines total)
%
% package coverage restored from the chapter's own preamble (2026-10-13): xy, xy-2cell

\usepackage[all]{xy}
\xyoption{2cell}
\UseAllTwocells

\newenvironment{reference}{\comment}{\endcomment}
\newenvironment{slogan}{\comment}{\endcomment}
\newenvironment{history}{\comment}{\endcomment}
\theoremstyle{plain}
\newtheorem{theorem}[subsection]{Theorem}
\newtheorem{proposition}[subsection]{Proposition}
\newtheorem{lemma}[subsection]{Lemma}
\theoremstyle{definition}
\newtheorem{definition}[subsection]{Definition}
\newtheorem{example}[subsection]{Example}
\newtheorem{exercise}[subsection]{Exercise}
\newtheorem{situation}[subsection]{Situation}
\theoremstyle{remark}
\newtheorem{remark}[subsection]{Remark}
\newtheorem{remarks}[subsection]{Remarks}
\numberwithin{equation}{subsection}
\def\lim{\mathop{\mathrm{lim}}\nolimits}
\def\colim{\mathop{\mathrm{colim}}\nolimits}
\def\Spec{\mathop{\mathrm{Spec}}}
\def\Hom{\mathop{\mathrm{Hom}}\nolimits}
\def\Ext{\mathop{\mathrm{Ext}}\nolimits}
\def\SheafHom{\mathop{\mathcal{H}\!\mathit{om}}\nolimits}
\def\SheafExt{\mathop{\mathcal{E}\!\mathit{xt}}\nolimits}
\def\Sch{\mathit{Sch}}
\def\Mor{\mathop{\mathrm{Mor}}\nolimits}
\def\Ob{\mathop{\mathrm{Ob}}\nolimits}
\def\Sh{\mathop{\mathit{Sh}}\nolimits}
\def\NL{\mathop{N\!L}\nolimits}
\def\CH{\mathop{\mathrm{CH}}\nolimits}
\def\proetale{{pro\text{-}\acute{e}tale}}
\def\etale{{\acute{e}tale}}
\def\QCoh{\mathit{QCoh}}
\def\Ker{\mathop{\mathrm{Ker}}}
\def\Im{\mathop{\mathrm{Im}}}
\def\Coker{\mathop{\mathrm{Coker}}}
\def\Coim{\mathop{\mathrm{Coim}}}
\DeclareMathSymbol{\boxtimes}{\mathbin}{AMSa}{"02}
\def\QCohstack{\mathcal{QC}\!\mathit{oh}}
\def\Cohstack{\mathcal{C}\!\mathit{oh}}
\def\Spacesstack{\mathcal{S}\!\mathit{paces}}
\def\Quotfunctor{\mathrm{Quot}}
\def\Hilbfunctor{\mathrm{Hilb}}
\def\Curvesstack{\mathcal{C}\!\mathit{urves}}
\def\Polarizedstack{\mathcal{P}\!\mathit{olarized}}
\def\Complexesstack{\mathcal{C}\!\mathit{omplexes}}
\def\Pic{\mathop{\mathrm{Pic}}\nolimits}
\def\Picardstack{\mathcal{P}\!\mathit{ic}}
\def\Picardfunctor{\mathrm{Pic}}
\def\Deformationcategory{\mathcal{D}\!\mathit{ef}}
\usepackage{cleveref}
\crefname{originallemma}{original lemma}{original lemmas}
\Crefname{originallemma}{Original lemma}{Original lemmas}
\makeatletter
\newlabel{algebraic-section-representable-by-algebraic-spaces}{{02ZV}{source}{}{sourcepointer-1}{}}
\newlabel{algebraic-section-representable-by-algebraic-spaces@cref}{{[originallemma][0][]02ZV}{[1][0][]source}{}{}{}}
\makeatother
\crefname{originallemma}{original lemma}{original lemmas}
\Crefname{originallemma}{Original lemma}{Original lemmas}
\makeatletter
\newlabel{algebraic-section-split}{{04SU}{source}{}{sourcepointer-2}{}}
\newlabel{algebraic-section-split@cref}{{[originallemma][0][]04SU}{[1][0][]source}{}{}{}}
\makeatother
\crefname{originallemma}{original lemma}{original lemmas}
\Crefname{originallemma}{Original lemma}{Original lemmas}
\makeatletter
\newlabel{sites-lemma-back-and-forth}{{00XT}{source}{}{sourcepointer-3}{}}
\newlabel{sites-lemma-back-and-forth@cref}{{[originallemma][0][]00XT}{[1][0][]source}{}{}{}}
\makeatother
\crefname{originallemma}{original lemma}{original lemmas}
\Crefname{originallemma}{Original lemma}{Original lemmas}
\makeatletter
\newlabel{sites-lemma-preserve-equalizers}{{00XS}{source}{}{sourcepointer-4}{}}
\newlabel{sites-lemma-preserve-equalizers@cref}{{[originallemma][0][]00XS}{[1][0][]source}{}{}{}}
\makeatother
\crefname{originallemma}{original lemma}{original lemmas}
\Crefname{originallemma}{Original lemma}{Original lemmas}
\makeatletter
\newlabel{sites-lemma-pullback-representable-presheaf}{{04D2}{source}{}{sourcepointer-5}{}}
\newlabel{sites-lemma-pullback-representable-presheaf@cref}{{[originallemma][0][]04D2}{[1][0][]source}{}{}{}}
\makeatother
\crefname{originallemma}{original lemma}{original lemmas}
\Crefname{originallemma}{Original lemma}{Original lemmas}
\makeatletter
\newlabel{sites-lemma-when-shriek}{{00XR}{source}{}{sourcepointer-6}{}}
\newlabel{sites-lemma-when-shriek@cref}{{[originallemma][0][]00XR}{[1][0][]source}{}{}{}}
\makeatother
\title{Stacks, Sheaves on Algebraic Stacks: lemma-compare}
\author{The Stacks Project Authors}
\date{stacks-project:stacks-sheaves.tex:local-snapshot-20261012; GNU FDL 1.2 or later}
\begin{document}
\maketitle
\noindent\textit{Editorial note.} Statement span is lines 1429-1452 of sources/stacks-algebraic-stacks-sheaves/pinned-original-source.tex; proof span is lines 1454-1505, the complete author proof (50 body lines).\par\medskip
\noindent\textit{Editorial note.} Navigation stubs: 6 cited-but-not-included labels are supplied as non-load-bearing pointers with their pinned tags where the snapshot tag index has them.\par\medskip
\noindent\textit{Editorial note.} extra\_packages is empty on purpose: the retained preamble already loads xy with options and a second declaration clashes in the assembler wrapper.\par\medskip
\noindent\textit{Editorial note.} Text is retained verbatim, including the author's own typos and placeholders.\par\medskip
\noindent\textit{Editorial note.} Queue index 8 of queue\_algebraic-stacks-sheaves.json (pb=2261, stmt 1429-1452, proof 1454-1505).\par\medskip
\noindent\textit{Allowed standard prerequisites (editorial).}\begin{itemize}
\item Explicit prerequisite boundary. The cited inputs are prerequisites; the argument in the retained span is not replaced by any cited result.
\end{itemize}\medskip
\ifdefined\hypertarget\hypertarget{sourcepointer-1}{}\fi
\noindent\textit{Source reference.} Original lemma 02ZV: sources/stacks-algebraic-stacks-sheaves/pinned-original-source.tex, source lines 1505--1505. Cited by the retained statement or proof but not part of the collected spans; supplied only as a navigation pointer. The argument does not depend on its content.\par\medskip
\ifdefined\hypertarget\hypertarget{sourcepointer-2}{}\fi
\noindent\textit{Source reference.} Original lemma 04SU: sources/stacks-algebraic-stacks-sheaves/pinned-original-source.tex, source lines 1505--1505. Cited by the retained statement or proof but not part of the collected spans; supplied only as a navigation pointer. The argument does not depend on its content.\par\medskip
\ifdefined\hypertarget\hypertarget{sourcepointer-3}{}\fi
\noindent\textit{Source reference.} Original lemma 00XT: sources/stacks-algebraic-stacks-sheaves/pinned-original-source.tex, source lines 1505--1505. Cited by the retained statement or proof but not part of the collected spans; supplied only as a navigation pointer. The argument does not depend on its content.\par\medskip
\ifdefined\hypertarget\hypertarget{sourcepointer-4}{}\fi
\noindent\textit{Source reference.} Original lemma 00XS: sources/stacks-algebraic-stacks-sheaves/pinned-original-source.tex, source lines 1505--1505. Cited by the retained statement or proof but not part of the collected spans; supplied only as a navigation pointer. The argument does not depend on its content.\par\medskip
\ifdefined\hypertarget\hypertarget{sourcepointer-5}{}\fi
\noindent\textit{Source reference.} Original lemma 04D2: sources/stacks-algebraic-stacks-sheaves/pinned-original-source.tex, source lines 1505--1505. Cited by the retained statement or proof but not part of the collected spans; supplied only as a navigation pointer. The argument does not depend on its content.\par\medskip
\ifdefined\hypertarget\hypertarget{sourcepointer-6}{}\fi
\noindent\textit{Source reference.} Original lemma 00XR: sources/stacks-algebraic-stacks-sheaves/pinned-original-source.tex, source lines 1505--1505. Cited by the retained statement or proof but not part of the collected spans; supplied only as a navigation pointer. The argument does not depend on its content.\par\medskip
\begin{lemma}
\label{lemma-compare}
Let $S$ be a scheme. Let $\mathcal{X} \to (\Sch/S)_{fppf}$ be a category
fibred in groupoids. Assume $\mathcal{X}$ is representable by an algebraic
space $F$. Then there exists a continuous and cocontinuous functor
$
F_\etale \to \mathcal{X}_\etale
$
which induces a morphism of ringed sites
$$
\pi_F :
(\mathcal{X}_\etale, \mathcal{O}_\mathcal{X})
\longrightarrow
(F_\etale, \mathcal{O}_F)
$$
and a morphism of ringed topoi
$$
i_F :
(\Sh(F_\etale), \mathcal{O}_F)
\longrightarrow
(\Sh(\mathcal{X}_\etale), \mathcal{O}_\mathcal{X})
$$
such that $\pi_F \circ i_F = \text{id}$. Moreover $\pi_{F, *} = i_F^{-1}$.
\end{lemma}

\begin{proof}
Choose an equivalence $j : \mathcal{S}_F \to \mathcal{X}$, see
Algebraic Stacks, Sections \ref{algebraic-section-split} and
\ref{algebraic-section-representable-by-algebraic-spaces}.
An object of $F_\etale$ is a scheme $U$ together with an
\'etale morphism $\varphi : U \to F$. Then $\varphi$ is an object
of $\mathcal{S}_F$ over $U$. Hence $j(\varphi)$ is an object of
$\mathcal{X}$ over $U$. In this way $j$ induces a functor
$u : F_\etale \to \mathcal{X}$. It is clear that
$u$ is continuous and cocontinuous for the \'etale topology on
$\mathcal{X}$. Since $j$ is an equivalence, the functor $u$ is fully
faithful. Also, fibre products and equalizers exist in $F_\etale$
and $u$ commutes with them because these are computed on the level
of underlying schemes in $F_\etale$. Thus
Sites, Lemmas \ref{sites-lemma-when-shriek},
\ref{sites-lemma-preserve-equalizers}, and
\ref{sites-lemma-back-and-forth}
apply. In particular $u$ defines a morphism of topoi
$i_F : \Sh(F_\etale) \to \Sh(\mathcal{X}_\etale)$
and there exists a left adjoint $i_{F, !}$ of $i_F^{-1}$ which commutes
with fibre products and equalizers.

\medskip\noindent
We claim that $i_{F, !}$ is exact. If this is true, then we can define
$\pi_F$ by the rules $\pi_F^{-1} = i_{F, !}$ and $\pi_{F, *} = i_F^{-1}$
and everything is clear. To prove the claim, note that we already know
that $i_{F, !}$
is right exact and preserves fibre products. Hence it suffices to show
that $i_{F, !}* = *$ where $*$ indicates the final object in the category
of sheaves of sets. Let $U$ be a scheme and let
$\varphi : U \to F$ be surjective and \'etale. Set $R = U \times_F U$.
Then
$$
\xymatrix{
h_R \ar@<1ex>[r] \ar@<-1ex>[r] & h_U \ar[r] & {*}
}
$$
is a coequalizer diagram in $\Sh(F_\etale)$. Using the
right exactness of $i_{F, !}$, using $i_{F, !} = (u_p\ )^\#$, and using
Sites, Lemma \ref{sites-lemma-pullback-representable-presheaf}
we see that
$$
\xymatrix{
h_{u(R)} \ar@<1ex>[r] \ar@<-1ex>[r] & h_{u(U)} \ar[r] & i_{F, !}{*}
}
$$
is a coequalizer diagram in $\Sh(\mathcal{X}_\etale)$. Using that
$j$ is an equivalence and that $F = U/R$ it follows that
the coequalizer in $\Sh(\mathcal{X}_\etale)$ of the
two maps $h_{u(R)} \to h_{u(U)}$ is $*$. We omit the proof that
these morphisms are compatible with structure sheaves.
\end{proof}

\end{document}
